\section{Diagnostic calibration}\label{app:diagnostics}
The independent curvature calibration uses two complete student trajectories
per task. The teaching policies are 7.5-second Short checkpoints. Calibration
students are unaudited proximal-target policies: 256-update fitting cycles,
full training-cost acceptance, teacher promotion on acceptance, and an
agreement-ratio coefficient schedule starting at 16. They use a 30-second
total allocation and separate teacher seed 9501. These policies supply calibration trajectories only.
Reverse differentiation evaluates each fixed teacher's value gradient at
the next state after its own action. The positive directional-remainder
coefficient is the larger of $2[C]_+/\|u-u_0\|^2$ for the tested student
action and proximal target, with a $10^{-20}$ denominator floor. Take its
95th percentile and multiply by 1.5 for the fixed empirical parameter.
This coefficient uses realized remainders and is not a prospective bound.

On these complete calibration trajectories, the maximum difference between
summed exact teacher advantages and the independently evaluated policy
cost difference is $3.94\times10^{-7}$. 

The perturbation experiment uses 150 task/teacher/quality/noise cases with
32 sampled actions each. It does not retrain students. Bound-violation
counts by condition (none, bias 0.5, Gaussian 0.5, bias 1.5, Gaussian 1.5)
are $(9,2,0,1,0)$ in mechanics and $(0,3,2,6,3)$ in reaction.
\begin{table}[htbp]\centering\small\setlength{\tabcolsep}{4pt}
\begin{tabular}{@{}lrrrr@{}}\toprule
Task & Noise & Naive FP/pred. & Margin FP/pred. & Coverage (\%)\\\midrule
Mechanical & none 0 & 2/480 & 2/480 & 100.0\\
Mechanical & bias 0.5 & 241/480 & 0/190 & 39.6\\
Mechanical & gaussian 0.5 & 296/480 & 0/98 & 20.4\\
Mechanical & bias 1.5 & 340/480 & 0/44 & 9.2\\
Mechanical & gaussian 1.5 & 413/480 & 0/24 & 5.0\\
Reaction & none 0 & 0/303 & 0/301 & 62.7\\
Reaction & bias 0.5 & 328/480 & 0/62 & 12.9\\
Reaction & gaussian 0.5 & 450/480 & 0/10 & 2.1\\
Reaction & bias 1.5 & 449/480 & 0/8 & 1.7\\
Reaction & gaussian 1.5 & 476/480 & 0/0 & 0.0\\
\bottomrule\end{tabular}
\caption{Independent local-action diagnostics over five teacher seeds and three checkpoint qualities: 480 action evaluations on 160 state/time points per task and noise condition, with states reused across teacher qualities. A fixed empirical curvature quantile comes from separate calibration trajectories. Ratios count false positive improvement predictions; zero predictions imply abstention. This calibrated margin is not a certified upper bound, and no student is retrained in this diagnostic.}
\end{table}


\section{Paired uncertainty}\label{app:uncertainty}
\begin{table}[htbp]\centering\small\setlength{\tabcolsep}{4pt}
\begin{tabular}{@{}lrrrr@{}}\toprule
Task & Reference & $\Delta$ cost & 95\% CI & Wins\\\midrule
Mechanical & Cloning & -0.000179 & [-0.016732, +0.016373] & 2/5\\
Mechanical & Cached H & -0.038344 & [-0.073294, -0.003394] & 5/5\\
Mechanical & Selected fixed & +0.001000 & [-0.007069, +0.009069] & 2/5\\
Mechanical & Adaptive & -0.003630 & [-0.023815, +0.016555] & 3/5\\
Reaction & Cloning & -0.015024 & [-0.071032, +0.040985] & 3/5\\
Reaction & Cached H & -0.092694 & [-0.161527, -0.023861] & 5/5\\
Reaction & Selected fixed & -0.106005 & [-0.255729, +0.043720] & 5/5\\
Reaction & Adaptive & -0.001017 & [-0.025796, +0.023762] & 3/5\\
\bottomrule\end{tabular}
\caption{Paired arithmetic cost differences across five teacher/student/data replications. Intervals use a two-sided $t$ calculation with four degrees of freedom; comparisons are exploratory and not multiplicity-adjusted.}
\end{table}

Each audit interval concerns the five paired teacher/student/data replications.
States within a trajectory and repeated actions under different perturbations
are not counted as independent replications. All selected policies and raw
per-state costs are supplied, including runs worse than their teacher.

\section{Reproducibility}\label{app:reproducibility}
The repository is
\url{https://github.com/yeoneung/teacher_student_hj}.
Its \href{https://github.com/yeoneung/teacher_student_hj/releases/tag/v1.0.0}{v1.0.0 release}
provides \texttt{research-artifacts-v1.zip}. Extracting this archive gives
\texttt{research/README.md} for environment and execution instructions,
\texttt{research/results/} for protocols, simulation outcomes and
checkpoints, and the analysis scripts in \texttt{research/}.
The release's \texttt{artifact-manifest.json} records asset and file hashes.
The reference archives \path{07_supporting_evidence.zip} and
\path{08_data_and_checkpoints.zip}, available in the same release,
are required by \texttt{verify\_evidence.py -{}-stage-inputs}.

The package supplies numerical-source and policy hashes, checkpoint
eligibility, raw costs, diagnostic arrays, environment versions and
entry points. Teacher construction, validation and test records remain
separate. Analysis re-evaluates stored policies without retraining.

The loss intervention includes 600 final policies and separate initial
and confirmation protocols. The confirmation plan hashes the initial
report and both numerical entry points before training. Checks verify
normal/radial decompositions and inactive-constraint parameter equality.
The minimal-information study supplies 480 policies, 80 caches with
actual query counts and a frozen protocol. Its analysis verifies
source/cache/checkpoint hashes and recovered coefficients. Separate
120-case algebraic checks cover sub-bound indistinguishability and
response error; these are not training replications.

The direct-query follow-up supplies 120 additional policies, all
re-evaluated, and scalar-count/timing records for 85 caches. Four
conditions reuse grid data for fidelity checks; the constructed
shared-actuation condition uses new seeds. The 2,304 approximate-feedback
checks are algebraic, not training replications. Across the distinct
studies the paper reports 1,306 fitted policies; their samples and
statistics are not pooled.
